\section{Conclusion}\label{sec:8}

This report analysed the VinaBeef requirements and identified three
driving characteristics: Reliability/Durability, Performance and
Security/Auditability. Availability is a supporting characteristic with
its own NFR, success criterion and scenario. A comparison of concrete
variants led to a service-based central backend with a gate core at
each farm. The design calls for:

\begin{itemize}
\item
  each gate read to commit at the farm within the two-second deadline,
  with no WAN or central dependency on its path;
\item
  committed facts to be relayed to the central services as events, at
  least once, with one business effect per intent;
\item
  each service to enforce role and scope on the resources it owns and
  commit the required audit evidence with each covered change.
\end{itemize}

The views, commit boundaries, role and custody policy, and
external-submission model form one consistent design, recorded in twelve
architecture decisions and supported by the design principles in
\Cref{sec:6}. In the M2 Health slice, the service integrated the
authorization policy and exercised retry, scope, and
source-evidence rollback against PostgreSQL.
\Cref{sec:m2-health-results} records the test results and remaining
mobile and identity integration checks. The other architecture
criteria in \Cref{sec:3} retain their own verification scope.
